\documentclass[10pt,a4paper]{amsart}
\usepackage[margin=19mm]{geometry}
\usepackage{amsmath,amssymb,mathtools}
\usepackage[scr=rsfs,cal=euler]{mathalfa}
\usepackage{xfrac}
\usepackage[unicode,hidelinks,bookmarks=false]{hyperref}
\newcommand{\ie}{i.e.\ }
\newcommand{\cf}{cf.\ }
\newcommand{\resp}{respectively\ }
\newcommand{\Subsection}{\subsection}
\providecommand{\nobreakdash}{\nobreak\hbox{-}}
\setlength{\emergencystretch}{2em}
\multlinegap0pt
\usepackage{bm}
\usepackage{bbm}
\usepackage{dsfont}
\usepackage{xspace}
\usepackage{cancel}
\usepackage{mathtools}
\usepackage{amsthm}
\makeatletter
\@ifundefined{defin}{\newtheorem{defin}{Defin}}{}
\@ifundefined{theorem}{\newtheorem{theorem}[defin]{Theorem}}{}
\makeatother
\title[n-th Tropical Nevanlinna Theory]{n-th Tropical Nevanlinna Theory}
\author{Risto Korhonen, Chengliang Tan}
\date{Source: arXiv:2602.03500v1 (CC BY 4.0)}
\begin{document}
\maketitle

\noindent\emph{Source and licence.} n-th Tropical Nevanlinna Theory. Version arXiv:2602.03500v1 (CC BY 4.0), distributed under the Creative Commons Attribution 4.0 International licence, as recorded in the arXiv licence metadata for this exact version and shown on the abstract page https://arxiv.org/abs/2602.03500v1. Full rights notice: licenses/arxiv-2602.03500-CC-BY-4.0.txt.\par\smallskip

\noindent\textit{Editorial scope.} Counted unit: the Theorem whose source label is \texttt{theorem5.8} in the article (source0.tex lines 1013--1023, source label \texttt{theorem5.8}), delivered together with the article's own complete original proof of it (source0.tex lines 1024--1169, one \texttt{proof} environment of 146 lines). No mathematics is written, repaired or re-derived: statement and proof are the article's own text line for line. Required context retained uncounted: none. Every definition, display and label of those lines is retained unchanged. Omitted from this package and not redistributed: 1--1012, 1170--1456. Nothing inside a retained excerpt is omitted or altered: every retained body is delivered exactly as the article's lines are written, so no omission receipt of any kind accompanies this package. Citations: the retained excerpts carry no citation command at all, so no bibliography entry of the article is transcribed and no reference is redistributed. Reference adaptation: none was needed and none was made; every reference inside the delivered unit resolves inside it, so no unresolved reference or citation is produced. Printed slips kept literal: no printed word, symbol, hypothesis or number of the counted statement or of its proof was altered. Layout and class: the article's own class is \texttt{article} with packages including graphicx, amsmath, amssymb, amsthm, tikz, pgfplots, booktabs, tabularx, multirow, authblk, cite; the wrapper is the lane's pinned \texttt{amsart} editorial wrapper at 10pt on A4, which restates the theorem-like environments the retained text needs and adds the article's own macro definitions that the retained text needs (none), with extra wrapper packages (bm, bbm, dsfont, xspace, cancel, mathtools). The delivered numbering is the unit's own. Assets: no retained span contains a figure environment, a graphics-inclusion command, a PDF-page inclusion, a table or a TikZ picture; no image file is copied into this package, the package declares no asset dependency and no image file is redistributed with it. Licence: version arXiv:2602.03500v1 of this article is distributed under the Creative Commons Attribution 4.0 International licence, as recorded in the arXiv licence metadata for this exact version and shown on the abstract page https://arxiv.org/abs/2602.03500v1. A full rights notice is delivered at licenses/arxiv-2602.03500-CC-BY-4.0.txt. Characters: every retained excerpt is pure ASCII, so the delivered engine, XeLaTeX with its default Latin Modern text font, reports no missing character.
\begin{theorem}[\textbf{Second main theorem for a homogeneous Fermat polynomial}]\label{theorem5.8}
    Let $m\in\mathbb{N}$ and $f:\mathbb{R}\rightarrow \mathbb{T}\mathbb{R}^{m}$ be a $1$-st tropical holomorphic curve with a reduced representation $\textbf{f}=(f_{0}, \cdots, f_{m})$ and $\mathcal{P}(x)=\sum_{i=0}^{m}\alpha_{i}x_{i}^{n}$ where $n\in\mathbb{N}$ and $\alpha_{i}>0$, $i=0, \cdots, m$. If
    \begin{eqnarray}\label{5f1}
        \limsup\limits_{r\rightarrow +\infty}\frac{r}{T_{f}(r)}=0,
    \end{eqnarray}
    then we have
    \begin{eqnarray}\label{5f2}
        \theta T_{f}^{n}(r)\leq \sum_{j=1}^{n}N^{(j)}\left(r, \frac{1_{0}}{\mathcal{P}\circ f}\oslash\right)+o(T_{f}^{n}(r))\leq \Theta T_{f}^{n}(r),
    \end{eqnarray}
    where $\theta=\min\limits_{0\leq i\leq m}\alpha_{i}$, $\Theta=2^{n-1}\sum_{i=0}^{m}\alpha_{i}$ and $r$ tends to infinity.
\end{theorem}

\begin{proof}
    For any linear function $g(x)=kx+b$, we can easily get 
    \begin{eqnarray*}
        \frac{(g^{n}(x))^{(j)}}{j!}=C_{n}^{j}g^{n-j}(x)k^{j},
    \end{eqnarray*}
    for all $j=0, \cdots, n$, $x\in\mathbb{R}$ where we assume $0^{0}=1$. Thus we can obtain
    \begin{eqnarray}\label{5a12}
        \omega_{f_{i}^{n}}^{(j)}(x)=C_{n}^{j}f_{i}^{n-j}(x)\left(sgn^{j+1}(x^{+})(f_{i}^{'}(x^{+}))^{j}-sgn^{j+1}(x^{-})(f_{i}^{'}(x^{-}))^{j}\right),
    \end{eqnarray}
    where $i=0, \cdots, m$ and $x\in\mathbb{R}$, from which we can see that
    \begin{eqnarray*}
        \{\text{the singularities of } f_{i}^{n}\}\subset\{\text{the singularities of } f_{i}\}\cup \{0\}.
    \end{eqnarray*}

    Since $f_{i}$ is $1$-st tropical entire, then it must satisfy one of the following two properties.

    (1) There exists a finite number $R_{i}>0$ such that $f_{i}^{'}(-r^{-})<0$ and $f_{i}^{'}(r^{+})>0$ when $r>R_{i}$;

    (2) $f_{i}^{'}(x^{\pm})\geq 0$ or $\leq 0$ holds for any $x\in\mathbb{R}$.

    If $f_{i}$ satisfies (1), then there exists a finite value $R_{i}^{'}\geq R_{i}$ such that $f_{i}(\pm r)>0$ for $r>R_{i}^{'}$, thus
    \begin{eqnarray*}
        M_{1, i}&:=&\max\left\{\max\limits_{x\in[-R_{i}^{'}, R_{i}^{'}]}\left\{|f_{i}(x)|\right\}, 1\right\}<+\infty,\\
        M_{2, i}&:=&\max\left\{\max\limits_{x\in[-R_{i}^{'}, R_{i}^{'}]}\left\{|f_{i}^{'}(x^{\pm})|\right\}, 1\right\}<+\infty.
    \end{eqnarray*}
    Then for any $x\in \mathbb{R}$, if $|x|> R_{i}^{'}$, we have $\omega_{f_{i}^{n}}^{(j)}(x)\geq 0$ from (\ref{5a12}), which implies that 
    \begin{eqnarray}\label{5a13}
      \sum_{j=1}^{n}N_{(-r, -R_{i}^{'})\cup (R_{i}^{'}, r)}^{(j)}(r, f_{i}^{n})=0.
    \end{eqnarray}
    
    If $|x|\leq R_{i}^{'}$, then we have $|\omega_{f_{i}^{n}}^{(j)}(x)|\leq C_{n}^{j}M_{1, i}^{n-j}2M_{2, i}^{j}\leq 2n!(M_{1, i}M_{2, i})^{n}=:M_{i}<\infty$ for all $j=1, \cdots, n$ from (\ref{5a12}). Suppose $\{x_{k}\}_{k=1}^{\nu_{i}}$ are the singularities of $f_{i}$ in $[-R_{i}^{'}, R_{i}^{'}]$, and since there is no finite accumulation point of $f_{i}$'s singularities, it follows that $\nu_{i}<+\infty$. Hence
    \begin{eqnarray}\label{5a14}
        \sum_{j=1}^{n}N_{[-R_{i}^{'}, R_{i}^{'}]}^{(j)}(r, f_{i}^{n})&\leq& \frac{1}{2}\sum_{j=1}^{n}\sum_{k=1}^{\nu_{i}}|\omega_{f_{i}^{n}}^{(j)}(x_{k})|(r-|x_{k}|)^{j}\nonumber\\
        &\leq& \frac{1}{2}n\nu_{i}M_{i}r^{n}=o(T_{f}^{n}(r)),
    \end{eqnarray}
    where we assume $r>\max\{R_{i}^{'}, 1\}$. The last equality follows by the assumption (\ref{5f1}).
    
    If $f_{i}$ satisfies (2), without loss of generality, we assume $f_{i}$ is non-constant and $f_{i}^{'}(x^{\pm})\geq 0$ for all $x\in\mathbb{R}$. For a small enough quantity $\epsilon$, we denote
    \begin{eqnarray*}
        I_{i}=:\{x:f_{i}(x+\epsilon)f_{i}(x-\epsilon)<0, x\in\mathbb{R}\}.
    \end{eqnarray*}
    Then there is at most one point in $I_{i}$. Let
    \begin{eqnarray*}
        x^{*}=\left\{\begin{array}{ll}
         x, & x\in I_{i},\\
         0, & I_{i}= \varnothing.
         \end{array}\right.
    \end{eqnarray*}
    Then $x^{*}$ is a finite number. If $x>\max\{x^{*}, 0\}$, then $f_{i}(x)>0$, and hence $\omega_{f_{i}^{n}}^{(j)}(x)\geq 0$ for $j=1, \cdots, 0$, which means that
    \begin{eqnarray}\label{5a15}
        \sum_{j=1}^{n}N_{(\max\{x^{*}, 0\}, r)}^{(j)}(r, f_{i}^{n})=0.
    \end{eqnarray}
    With the same calculation of (\ref{5a14}), we can also get 
    \begin{eqnarray}\label{5a16}
        \sum_{j=1}^{n}N_{[\min\{x^{*}, 0\}, \max\{x^{*}, 0\}]}^{(j)}(r, f_{i}^{n})=o(T_{f}^{n}(r)).
    \end{eqnarray}
    
    If $x<\min\{x^{*}, 0\}$, we assume $f_{i, x^{*-}}(x)=k_{i}^{*}x+b_{i}^{*}$ where $k_{i}^{*}, b_{i}^{*}\in\mathbb{R}$ are finite values and $k_{i}^{*}\geq 0$. Since $f_{i}$ is non-decreasing and convex, then $|f(x)-b_{i}^{*}|\leq |k_{i}^{*}x|$ for $x<\min\{x^{*}, 0\}$, which leads to
    \begin{eqnarray}\label{5a17}
        |f_{i}(x)|\leq k_{i}^{*}|x|+O(1),
    \end{eqnarray}
    where $x<\min\{x^{*}, 0\}$. For any $r>|x^{*}|$ and then $f(r)\geq 0$, we denote by $\{x_{-l}\}_{l=1}^{\mu_{i}}$ the roots of $f_{i}$ in $(-r, \min\{x^{*}, 0\})$ and define a new $1$-st tropical entire function
    \begin{eqnarray*}
        g_{i}(x):=\left\{\begin{array}{ll}
         k_{i}^{*}x+b_{i}^{*}, & x> \min\{x^{*}, 0\},\\
         f_{i}(x), & x\leq \min\{x^{*}, 0\}.
         \end{array}\right.
    \end{eqnarray*}
    Then $m(r, \pm g_{i})\leq k_{i}^{*}r+O(1)$ and $\{x_{-l}\}_{l=1}^{\mu_{i}}$ are all $g_{i}$'s roots in $(-r, r)$. By following the $1$-st Jensen formula, we have
    \begin{eqnarray*}
        N(r, 1_{0}\oslash g_{i})&=&\frac{1}{2}\sum_{l=1}^{\mu_{i}}(f_{i}^{'}(x_{-l}^{+})-f_{i}^{'}(x_{-l}^{-}))(r-|x_{-l}|)\\
        &=&T(r, 1_{0}\oslash g_{i})-m(r, 1_{0}\oslash g_{i})\\
        %&=&T(r, 1_{0}\oslash g_{i})-\frac{1}{2}(-g_{i}(-r))\\
        &\leq&T(r, g_{i})+m(r, 1_{0}\oslash g_{i})+O(1)\\
        &=&m(r, g_{i})+m(r, 1_{0}\oslash g_{i})+O(1)\\
        &\leq&2k_{i}^{*}r+O(1).
    \end{eqnarray*}
    Therefore, from (\ref{5a12}), we have 
    \begin{eqnarray}\label{5a18}
        &&\sum_{j=1}^{n}N_{(-r, \min\{x^{*}, 0\})}^{(j)}(r, f_{i}^{n})\nonumber\\
        &\leq&\sum_{j=1}^{n}\sum_{l=1}^{\mu_{i}}n!\left(k_{i}^{*}r+O(1)\right)^{n-j}\left((f_{i}^{'}(x_{-l}^{+}))^{j}-(f_{i}^{'}(x_{-l}^{-}))^{j}\right)(r-|x_{-l}|)^{j}\nonumber\\
        &\leq&\sum_{j=1}^{n}n!\left(k_{i}^{*}r+O(1)\right)^{n-j}\sum_{l=1}^{\mu_{i}}\left((f_{i}^{'}(x_{-l}^{+})-f_{i}^{'}(x_{-l}^{-}))(r-|x_{-l}|)n^{2}(k^{*}r)^{j-1}\right)\nonumber\\
        &\leq&\sum_{j=1}^{n}(n+2)!\left(k_{i}^{*}r+O(1)\right)^{n-1}\left(\sum_{l=1}^{\mu_{i}}(f_{i}^{'}(x_{-l}^{+})-f_{i}^{'}(x_{-l}^{-}))(r-|x_{-l}|)\right)\nonumber\\
        &\leq&\sum_{j=1}^{n}(n+2)!\left(k_{i}^{*}r+O(1)\right)^{n-1}\left(4k_{i}^{*}r+O(1)\right)\nonumber\\
        &\leq&4(n+3)!\left(k_{i}^{*}r+O(1)\right)^{n}\nonumber\\
        &=&o(T_{f}^{n}(r))
    \end{eqnarray}
    where we used $0\leq f_{i}^{'}(x_{-l}^{\pm})\leq k^{*}$ for $l=1, \cdots, \mu_{i}$ and also used a basic inequality $a^{j}-b^{j}=(a-b)(\sum_{s=0}^{j-1}a^{s}b^{j-s-1})\leq n(a-b)a^{j-1}$ for $0<b\leq a$ and $j=1, \cdots, n$ in the second inequality. We can also do a similar analysis in the case where $f_{i}^{'}(x^{\pm})\leq 0$ for all $x\in\mathbb{R}$.

    By combining equations (\ref{5a13}), (\ref{5a14}), (\ref{5a15}), (\ref{5a16}) and (\ref{5a18}), we obtain
    \begin{eqnarray}\label{5a19}
        \sum_{j=1}^{n}N^{(j)}(r, f_{i}^{n})=o(T_{f}^{n}(r)).
    \end{eqnarray}
    
    For any $n$-th tropical meromorphic functions $h_{1}$, $h_{2}$ and any positive constant $\alpha$, we can easily get that
    \begin{eqnarray*}
        \sum_{j=1}^{n}N^{(j)}(r, h_{1}+h_{2})&\leq& \sum_{j=1}^{n}N^{(j)}(r, h_{1})+\sum_{j=1}^{n}N^{(j)}(r, h_{2})\\
        \sum_{j=1}^{n}N^{(j)}(r, \alpha h_{1})&=&\alpha \sum_{j=1}^{n}N^{(j)}(r, h_{1}).
    \end{eqnarray*}
    Hence
    \begin{eqnarray*}
        \sum_{j=1}^{n}N^{(j)}(r, \mathcal{P}\circ f)\leq \sum_{i=0}^{m}\alpha_{i}\sum_{j=1}^{n}N^{(j)}(r, f_{i}^{n})=o(T_{f}^{n}(r)).
    \end{eqnarray*}

    Since $\mathcal{P}\circ f$ is a $n$-th tropical meromorphic function, then by applying the $n$-th Jensen formula, we have
    \begin{eqnarray}\label{Pn}
        &&\sum_{j=1}^{n}N^{(j)}\left(r, \frac{1_{0}}{\mathcal{P}\circ f}\oslash\right)\nonumber\\
        &=&m(r, \mathcal{P}\circ f)-m\left(r, \frac{1_{0}}{\mathcal{P}\circ f}\oslash\right)+\sum_{j=1}^{n}N^{(j)}\left(r, \mathcal{P}\circ f\right)+O(1)\nonumber\\
        &=&\frac{1}{2}\sum_{\delta=\pm 1}\mathcal{P}\circ f(\delta r)+o(T_{f}^{n}(r)).
    \end{eqnarray}

From (\ref{5a17}), when $r$ is sufficient large, $f_{i}$ satisfies one of the following 3 properties

(1) $f_{i}(\pm r)\geq 0$;

(2) $f_{i}(r)\geq 0$, $f_{i}(-r)<0$ and $|f_{i}(-r)|\leq kr+O(1)$ where $k$ is a finite positive constant;

(3) $f_{i}(-r)\geq 0$, $f_{i}(r)<0$ and $|f_{i}(r)|\leq kr+O(1)$ where $k$ is a finite positive constant.

    Let $F(x)=\max\limits_{0\leq i\leq m}\{f_{i}(x)\}$, then $F(\pm r)$ can not both be negative simultaneously when $r$ is large enough, because all $f_{i}$ are tropical entire and $f$ is not a constant holomorphic curve. Now we claim that
    \begin{eqnarray}\label{5a11}
        \left(\frac{F(r)+F(-r)}{2}\right)^{n}\leq \frac{F^{n}(r)+F^{n}(-r)}{2}+o(T_{f}^{n}(r))\leq \frac{(F(r)+F(-r))^{n}}{2}
    \end{eqnarray}
    as $r$ tends to infinity. 
    
    We first suppose $F(\pm r)>0$ for all large enough $r$. If $F(-\delta r)=o(F(\delta r))$ for some $\delta=\pm 1$, then (\ref{5a11}) becomes $(F(\delta r)/2)^{n}\leq F^{n}(\delta r)/2\leq F^{n}(\delta r)/2$ which is valid. If $F(-\delta r)= a F(\delta r)(1+o(1))$ for some constants $a\geq 0$ and $\delta=\pm 1$, without loss of generality, we assume $a\geq 1$, otherwise we can consider $F(\delta r)= \frac{1}{a}F(-\delta r)(1+o(1))$. Then (\ref{5a11}) can be rewritten as $\left(\frac{a+1}{2}\right)^{n}F^{n}(\delta r)\leq \frac{1+a^{n}}{2}F^{n}(\delta r)\leq \frac{(1+a)^{n}}{2}F^{n}(\delta r)$ which is valid as well; the first inequality follows by considering the derivative of $\eta(t)=(1+(2t-1)^{n})/2-t^{n}$ where $t=\frac{1+a}{2}\geq 1$ and $n\in\mathbb{N}$.

    If $F(-\delta r)<0$ for large enough $r$ and $\delta=1$ or $-1$, then from above properties of $f_{i}$, we know that $|F(-\delta r)|=o(T_{f}(r))$ and $F(\delta r)=2T_{f}(r)+o(T_{f}(r))$, which implies $F(-\delta r)=o(F(\delta r))$. Then the claim follows by a similar analysis as above.
    
    Hence, we complete the proof of the claim. By applying (\ref{5a11}), it follows, on one hand, that
    \begin{eqnarray}\label{Pn1}
        \frac{1}{2}\sum_{\delta=\pm 1}\mathcal{P}\circ f(\delta r)&=&\frac{1}{2}\sum_{\delta=\pm 1}\sum_{i=0}^{m}\alpha_{i}f_{i}^{n}(\delta r)\nonumber\\
        &\leq&\frac{1}{2}\sum_{\delta=\pm 1}\left(\sum_{i=0}^{m}\alpha_{i}\right)F^{n}(\delta r)+o(T_{f}^{n}(r))\nonumber\\
        %&\leq&\left(\sum_{i=0}^{m}\alpha_{i}\right)\frac{(F(r)+F(-r))^{n}}{2}+o(T_{f}^{n}(r))\nonumber\\
        &\leq&2^{n-1}\left(\sum_{i=0}^{m}\alpha_{i}\right)\left(\frac{F(r)+F(-r)}{2}\right)^{n}+o(T_{f}^{n}(r))\nonumber\\
        &=&\Theta T_{f}^{n}(r)+o(T_{f}^{n}(r)),
    \end{eqnarray}
    and, on the other hand, 
    \begin{eqnarray}\label{Pn2}
        \frac{1}{2}\sum_{\delta=\pm 1}\mathcal{P}\circ f(\delta r)&=&\frac{1}{2}\sum_{\delta=\pm 1}\sum_{i=0}^{m}\alpha_{i}f_{i}^{n}(\delta r)\nonumber\\
        &\geq&\frac{1}{2}\theta\sum_{\delta=\pm 1}F^{n}(\delta r)+o(T_{f}^{n}(r))\nonumber\\
        &\geq& \theta\left(\frac{F(r)+F(-r)}{2}\right)^{n}+o(T_{f}^{n}(r))\nonumber\\
        &=&\theta T_{f}^{n}(r)+o(T_{f}^{n}(r)).
    \end{eqnarray}
    Thus the assertion follows by combining (\ref{Pn}), (\ref{Pn1}) and (\ref{Pn2}).
\end{proof}

\end{document}
